% =============================================================
% Distance-Aware Loss for Phonologically-Graded G2P in PT-BR
% IEEE SLT 2026 — Double-blind submission (6 pages + 2 refs)
% =============================================================
%
% COMPILATION:
%   xelatex main && bibtex main && xelatex main && xelatex main
%   OR
%   pdflatex main && bibtex main && pdflatex main && pdflatex main
% =============================================================

\documentclass[conference]{IEEEtran}
\IEEEoverridecommandlockouts

% ---- Encoding & fonts ----------------------------------------
\usepackage[utf8]{inputenc}
\usepackage[T1]{fontenc}
\usepackage{tipa}

% ---- Math ----------------------------------------------------
\usepackage{amsmath}
\usepackage{amssymb}

% ---- Tables --------------------------------------------------
\usepackage{booktabs}
\usepackage{multirow}

% ---- References ----------------------------------------------
\usepackage{cite}

% ---- Macros --------------------------------------------------
\newcommand{\dpp}{d_{\mathrm{PP}}}        % PanPhon distance
\newcommand{\LCE}{L_{\mathrm{CE}}}        % CE loss
\newcommand{\lam}{\lambda}                % lambda shorthand

% =============================================================
\begin{document}
% =============================================================

\title{Distance-Aware Loss for Phonologically-Graded\\
       Grapheme-to-Phoneme Conversion in Brazilian Portuguese}

\author{%
  \IEEEauthorblockN{[ANONYMOUS]}
  \IEEEauthorblockA{[ANONYMOUS]}
}

\maketitle

% ---- Abstract -----------------------------------------------
\begin{abstract}
Standard grapheme-to-phoneme (G2P) training with cross-entropy
treats all phoneme substitutions equally, regardless of articulatory
distance.
We propose Distance-Aware (DA) Loss, which penalizes substitutions
proportionally to PanPhon articulatory distance between predicted and
target phonemes, weighted by the model's prediction confidence.
Applied to Brazilian Portuguese with a BiLSTM encoder-decoder,
DA~Loss systematically redistributes errors toward phonologically
closer targets: Class~D (catastrophic) substitutions decrease
19\,\% relative vs.\ a CE~baseline, while Class~B (near-miss)
errors increase proportionally.

Our system achieves PER~0.48\,\% and WER~5.33\,\% on a stratified
28{,}782-word test set---57$\times$ larger than comparable PT-BR
evaluations---with a Wilson 95\,\% CI of $\pm$0.03\,pp.
Both Wilson CIs do not overlap with the closest reference system
(upper bound 0.51\,\% vs.\ lower bound 0.56\,\%)---a statistically
significant difference at 95\,\% confidence.
Evaluation on 31~out-of-vocabulary words yields 100\,\% accuracy
on genuine novel PT-BR words, consistent with phonological rule
generalization beyond memorization.
A factorial analysis of model capacity, syllable separators, and
DA~Loss reveals that the interaction between these factors---not
capacity alone---determines performance.
\end{abstract}

\begin{IEEEkeywords}
grapheme-to-phoneme, distance-aware loss, Brazilian Portuguese,
BiLSTM, phonological features, articulatory distance
\end{IEEEkeywords}

% =============================================================
\section{Introduction}
% =============================================================

Grapheme-to-phoneme (G2P) conversion is a core component of
text-to-speech synthesis, automatic speech recognition, and
multilingual NLP pipelines.
For Brazilian Portuguese (PT-BR), G2P presents well-documented
challenges: graphemic ambiguity (grapheme ``c'' maps to /k/ in
\textit{cama} but /s/ in \textit{cena}; ``r'' maps to
/\textipa{\*r}/ in syllable onset but /x/ in word-final
coda~\cite{barbosa2004brazilian}), vowel neutralization in
unstressed positions (/e/$\leftrightarrow$/\textipa{E}/ and
/o/$\leftrightarrow$/\textipa{O}/ merge), and position-dependent
coda realization (/x/ in final position; /\textipa{G}/ before
voiced consonants).
Prior work on PT-BR G2P has addressed these with decision trees,
n-gram models, and neural seq2seq
architectures~\cite{neto2006brazilian}.

Standard models trained with cross-entropy (CE) treat all phoneme
errors equally: predicting /\textipa{E}/ when the target is /e/---a
near-miss differing in one articulatory feature---incurs the same
gradient penalty as predicting /k/ for /a/, spanning eight features.
This phonological blindness distorts the training signal: the model
learns \emph{that} it erred but not \emph{how severely}.

We address this with \textbf{Distance-Aware (DA) Loss}, which adds a
training signal proportional to the articulatory distance between
predicted and target phonemes, weighted by prediction confidence.
We apply this to Brazilian Portuguese using a BiLSTM encoder-decoder
with Bahdanau attention~\cite{bahdanau2014neural}---deliberately
chosen to isolate the contribution of the loss function from
architectural novelty.

\textbf{Comparison with LatPhon~\cite{chary2025latphon}.}
LatPhon is a 4-layer multilingual Transformer (7.5M parameters, RoPE)
reporting PER~0.86\,\% (Wilson 95\,\% CI [0.56\,\%, 1.16\,\%]) on
$\approx$500~PT-BR words from the same IPA dictionary used here.
Our system achieves PER~0.48\,\% (CI~[0.46\,\%, 0.51\,\%]) on
28{,}782~words.
Both systems report Wilson CIs; the intervals do not overlap---our
upper bound (0.51\,\%) falls below LatPhon's lower bound
(0.56\,\%), a statistically significant difference at 95\,\%
confidence.
The 57$\times$ difference in test set size confers 10$\times$ more
precise CIs to our evaluation ($\pm$0.03\,pp vs.\
$\pm$0.30\,pp).

\textbf{Contributions:}
\begin{enumerate}
  \item DA~Loss: a phonologically-graded training objective combining
        PanPhon articulatory distance with prediction confidence
        (Sec.~\ref{sec:daloss})
  \item Large-scale stratified evaluation with empirical evidence of
        split bias impact (Sec.~\ref{sec:data})
  \item Factorial analysis of capacity $\times$ separators $\times$
        DA~Loss interaction (Sec.~\ref{sec:factorial})
  \item Error quality taxonomy: Class~A--D classification revealing
        systematic redistribution of catastrophic errors
        (Sec.~\ref{sec:analysis})
  \item OOV generalization evaluation across 6~diagnostic categories
        (Sec.~\ref{sec:oov})
\end{enumerate}

% =============================================================
\section{Related Work}
\label{sec:related}
% =============================================================

Neural G2P systems have evolved from early RNN approaches to modern
Transformer architectures.
Rao et al.~\cite{rao2015g2p} demonstrated that LSTM-based
encoder-decoders outperform n-gram joint-sequence
models~\cite{bisani2008joint} on English G2P.
For Portuguese, Neto et al.~\cite{neto2006brazilian} established
early baselines with decision trees and HMMs.
More recently, LatPhon~\cite{chary2025latphon} applied a
multilingual Transformer with RoPE to 6~languages including PT-BR.

\textbf{Phonological distance in training.}
Distance-weighted objectives have been explored in speech
recognition and synthesis, where confusing acoustically similar
phonemes is less costly than distant substitutions.
Our work differs in three aspects: (1)~we operate at the IPA
grapheme-to-phoneme level rather than acoustic features;
(2)~we use PanPhon articulatory distance~\cite{mortensen2016panphon}
as a principled, language-independent distance metric;
(3)~we weight the distance penalty by prediction confidence,
creating a composite signal that maximally penalizes errors that are
both confident and phonologically distant.

\textbf{Evaluation methodology.}
G2P evaluations typically use small test sets (300--1{,}000 words)
with Wald confidence intervals.
We evaluate on 28{,}782 stratified words with Wilson
CIs~\cite{wilson1927probable}, providing 10$\times$ more precise
uncertainty estimates.

% =============================================================
\section{Data and Evaluation Protocol}
\label{sec:data}
% =============================================================

\subsection{Corpus}

The training corpus consists of \textbf{95{,}937 grapheme--IPA pairs}
from a Brazilian Portuguese phonetic dictionary.
The input charset covers a--z (excluding k, w, y, absent from the
dictionary) plus Portuguese diacritics.
Words containing k, w, or~y are character-OOV and mapped to
$\langle$UNK$\rangle$.
Prior to training, 10{,}252 entries were corrected for ASCII-g
(U+0067) vs.\ IPA-\textipa{g} (U+0261) conflation, required for
correct PanPhon feature lookup.

\subsection{Stratified Split}

The corpus is divided \textbf{60/10/30} (train/val/test) using
stratified sampling over three phonological variables: stress type
(oxy\-tone/paroxytone/proparoxytone), syllable count bin
(1, 2, 3, 4, 5+), and character length bin
($\leq$4, 5--7, 8--10, 11+), yielding approximately 48~strata.
A purely random split risks concentrating easy words in the test set,
inflating metrics.
Split quality: $\chi^2\!=\!0.95$ ($p\!=\!0.678$),
Cram\'{e}r $V\!=\!0.0007$---no statistically significant
distributional difference across subsets.

\textbf{Empirical evidence of split bias.}
An unstratified 70/10/20 split (Exp0, same 4.3M architecture, CE
loss) achieved PER~1.12\,\%, while the stratified 60/10/30 split
(Exp1) achieved 0.66\,\%---a 41\,\% PER reduction attributable
entirely to evaluation protocol, not model
improvement~\cite{kohavi1995crossvalidation}.
Without stratification, the random partition can concentrate difficult
words (proparoxytones, long words) in training and easy words in
testing.
This motivated stratification in all subsequent experiments and the
deliberate choice of a large 30\,\% test set over maximizing
training data.

\subsection{Evaluation Metrics}

\textbf{PER} (Phoneme Error Rate): Levenshtein distance over phoneme
sequences, normalized by reference length~\cite{bisani2008joint}:
\begin{equation}
  \mathrm{PER} = \frac{\displaystyle\sum_i \mathrm{edit}(\hat{y}_i,\,
                 y_i)}{\displaystyle\sum_i |y_i|} \times 100\%
  \label{eq:per}
\end{equation}

\textbf{WER} (G2P Word Error Rate): fraction of words with any
phoneme error (exact-match).
\textbf{Wilson 95\,\% CI}~\cite{wilson1927probable,brown2001interval}:
used throughout in place of Wald intervals, which underestimate
uncertainty near $\hat{p}\!\to\!0$.
For our test set ($\approx$181K reference phonemes), the Wilson CI on
PER is $\pm$0.03\,pp---10$\times$ more precise than
$\pm$0.30\,pp for $\approx$500-word evaluations.

% =============================================================
\section{Architecture}
% =============================================================

We use a \textbf{BiLSTM encoder-decoder with Bahdanau
attention}~\cite{bahdanau2014neural}, established for supervised
G2P~\cite{rao2015g2p}, and chosen to isolate the loss contribution
from architectural novelty.

\textbf{Encoder}: 2-layer BiLSTM over grapheme embeddings.
Each position~$t$ yields $h_t\!=\![\overrightarrow{h}_t;\,
\overleftarrow{h}_t]$, providing full bidirectional context---critical
for resolving graphemic ambiguity.

\textbf{Attention}: at each decoder step~$t$:
\begin{align}
  e_{t,j} &= v^\top \tanh(W_h h_j + W_s s_{t-1}), \nonumber \\
  \alpha_{t,j} &= \mathrm{softmax}(e_{t,j}), \quad
  c_t = \textstyle\sum_j \alpha_{t,j}\, h_j
  \label{eq:attention}
\end{align}

\textbf{Decoder}: 2-layer LSTM with teacher forcing during training;
autoregressive at inference.

\textbf{Phoneme embeddings.}
All main experiments use learned embeddings (random Glorot init).
We additionally tested PanPhon-initialized embeddings (24-D
articulatory features projected to 128-D), providing a geometric
prior where similar phonemes start close in embedding space.
With CE loss, PanPhon init produces qualitatively better errors
(lower weighted PER) but does not improve raw PER, as CE does not
reinforce the articulatory structure.
DA~Loss operates independently of embedding initialization: it uses
an external lookup table of PanPhon distances, not the embedding
geometry.
The two mechanisms are orthogonal and could potentially be combined;
this remains unexplored.

Three capacity configurations were evaluated
(Table~\ref{tab:configs}).

\begin{table}[t]
  \caption{Architecture configurations}
  \label{tab:configs}
  \centering\small
  \begin{tabular}{lccc}
    \toprule
    Config  & Embedding & Hidden & Parameters \\
    \midrule
    Small   & 128-D & 256-D & 4.3M \\
    Medium  & 192-D & 384-D & 9.7M \\
    Large   & 256-D & 512-D & 17.2M \\
    \bottomrule
  \end{tabular}
\end{table}

\subsection{Training Protocol}

All models are trained with Adam ($\beta_1\!=\!0.9$,
$\beta_2\!=\!0.999$), learning rate $10^{-3}$, batch size 96,
and early stopping with patience 10 on validation loss.
Teacher forcing ratio is 1.0 during training.
Batch size 96 was chosen empirically: with 34 phonological strata
in the corpus, $\text{batch\_size} \approx 3 \times S$ ensures
$\geq$3 samples per stratum per batch, reducing loss curve
variance by $\sim$50\,\% relative to batch size 32.
Models converge in 80--120 epochs; superior configurations
consistently converge in fewer epochs with lower validation loss,
supporting early stopping as a reliable quality indicator.
All experiments use fixed random seed (42) and deterministic
operations where available.

% =============================================================
\section{Distance-Aware Loss}
\label{sec:daloss}
% =============================================================

\subsection{Formulation}

\begin{equation}
  \boxed{L = \LCE + \lam \cdot \dpp(\hat{y},\, y) \cdot p(\hat{y})}
  \label{eq:daloss}
\end{equation}

\begin{itemize}
  \item $\LCE = -\log p_\mathrm{correct}$: standard cross-entropy
  \item $\dpp(\hat{y}, y) \in [0,1]$: normalized Euclidean distance
        in PanPhon's 24-dimensional articulatory feature space
        encoding voicing, nasality, place, and manner of
        articulation~\cite{mortensen2016panphon}.
        Example values: p$\leftrightarrow$b: 0.04 (voicing only);
        s$\leftrightarrow$\textipa{S}: 0.15 (place);
        a$\leftrightarrow$k: 0.90 (vowel vs.\ velar stop)
  \item $p(\hat{y}) \in (0,1]$: softmax probability of the
        \emph{predicted} (argmax) phoneme---independent of
        $p_\mathrm{correct}$
  \item $\lam$: coupling weight ($\lam\!=\!0.20$ throughout)
\end{itemize}

\textbf{Design rationale.}
CE monitors $p_\mathrm{correct}$; DA monitors $p(\hat{y})$.
These are independent: when the model is confidently wrong,
$p_\mathrm{correct}$ is low while $p(\hat{y})$ is high.
The confidence factor scales the articulatory penalty with the
model's certainty---maximally penalizing errors that are both
confident and phonologically distant.

\textbf{Bounded signal.}
DA~Loss is upper-bounded by $\lam\!\times\!1.0\!\times\!1.0=0.20$,
while CE can reach $\approx$16 in early training.
DA is effective primarily in the learning transition zone
(CE~0.3--1.5), where the model is actively resolving phoneme
ambiguities.

\textbf{Numerical example.}
Consider two errors with identical CE ($p_\mathrm{correct}\!=\!0.40$,
CE$\,=\,0.916$):
a near-miss e$\to$\textipa{E} ($d\!=\!0.10$, $p(\hat{y})\!=\!0.45$)
adds DA$\,=\,0.009$, total $L\!=\!0.925$;
a catastrophic a$\to$k ($d\!=\!0.90$, $p(\hat{y})\!=\!0.45$) adds
DA$\,=\,0.081$, total $L\!=\!0.997$.
CE alone is identical (0.916) in both cases.
DA provides a 9$\times$ difference in phonological penalty, steering
the model to ``break ties toward the closer candidate'' over training
epochs.

\textbf{Novelty scope.}
Distance-weighted losses have been explored in ASR and TTS.
Our contribution is the specific compositional formulation for
IPA-level G2P: PanPhon distance as gradient signal, prediction
confidence as weighting factor, and controlled coupling via~$\lam$
preserving CE as primary learning axis.

\subsection{Structural Token Override}

PanPhon assigns zero vectors to non-phonemic tokens: syllable
boundary ``\textbf{.}'' and stress marker ``\textbf{'}'' (U+02C8).
Without correction, $d(., \text{'}) = 0$---DA provides no gradient
signal for their confusion.
We apply a post-normalization override: distance~1.0 between any
structural token and any other symbol.
\emph{The override must be applied after Euclidean normalization};
applying it before yields $d\!\approx\!0.25$, not the intended
maximum.

\subsection{$\lambda$ Sweep}

Table~\ref{tab:lambda} sweeps $\lam$ with fixed architecture (4.3M,
no syllable separators).
The inverted-U curve motivates $\lam\!=\!0.20$: too low leaves CE
undifferentiated; too high competes with CE.

\begin{table}[t]
  \caption{$\lambda$ sweep, 4.3M architecture, no separators}
  \label{tab:lambda}
  \centering\small
  \begin{tabular}{cccc}
    \toprule
    $\lambda$ & PER & WER & Behavior \\
    \midrule
    0.05 & 0.62\,\% & 5.36\,\% & Signal too weak \\
    0.10 & 0.63\,\% & 5.35\,\% & Moderate improvement \\
    \textbf{0.20} & \textbf{0.60\,\%} & \textbf{5.14\,\%} & \textbf{Optimal} \\
    0.50 & 0.65\,\% & 5.57\,\% & Over-penalization \\
    \bottomrule
  \end{tabular}
\end{table}

% =============================================================
\section{Experiments}
\label{sec:results}
% =============================================================

\subsection{Main Results}

Table~\ref{tab:results} presents the main experimental results.
Stratified evaluation across 28{,}782 test words yields Wilson CI
$\pm$0.03\,pp---57$\times$ larger than comparable PT-BR
evaluations.

\begin{table}[t]
  \caption{Main results. Sep.\ = syllable boundary tokens.
           Dist.\ = structural distance override.
           Bold = recommended configurations.}
  \label{tab:results}
  \centering\small
  \begin{tabular}{lccccc}
    \toprule
    System & Params & Loss & Sep. & PER & WER \\
    \midrule
    CE baseline         & 4.3M  & CE      & ---        & 0.66\,\% & 5.65\,\% \\
    DA $\lam\!=\!0.2$   & 4.3M  & DA      & ---        & 0.60\,\% & 5.14\,\% \\
    \textbf{DA} $\boldsymbol{\lam\!=\!0.2}$ & \textbf{9.7M} & \textbf{DA} & --- & \textbf{0.58\,\%} & \textbf{4.96\,\%} \\
    CE + sep.           & 9.7M  & CE      & \checkmark & 0.52\,\% & 5.79\,\% \\
    DA + sep.           & 9.7M  & DA      & \checkmark & 0.53\,\% & 5.73\,\% \\
    \textbf{DA+dist}    & \textbf{17.2M} & \textbf{DA+dist} & \checkmark & \textbf{0.48\,\%} & \textbf{5.33\,\%} \\
    \midrule
    LatPhon~\cite{chary2025latphon} & 7.5M & CE & --- & 0.86\,\% & n/a \\
    \bottomrule
  \end{tabular}
\end{table}

\textbf{Separator trade-off.}
Syllable boundary tokens improve PER ($-17$\,\% to $-20$\,\%) at
the cost of WER ($+6$\,\% to $+8$\,\%).
Each misplaced separator counts as a full word error under
exact-match WER---a structural trade-off independent of architecture
or loss.
\textit{Recommendation}: use the 9.7M DA model (no separators)
for WER-sensitive tasks; use the 17.2M DA+dist model for
PER-sensitive tasks (TTS, forced alignment).

\subsection{Factorial Analysis}
\label{sec:factorial}

A 2$\times$2 factorial (separators $\times$ DA~Loss) at 9.7M
parameters isolates individual factor contributions
(Table~\ref{tab:factorial}).

\begin{table}[t]
  \caption{2$\times$2 factorial: separators $\times$ DA Loss,
           9.7M parameters, 60/10/30 split. Values are PER (\%).}
  \label{tab:factorial}
  \centering\small
  \begin{tabular}{lcc}
    \toprule
    & CE & DA $\lam\!=\!0.2$ \\
    \midrule
    No separators   & 0.63\,\% & 0.58\,\% \\
    With separators & 0.52\,\% & 0.53\,\% \\
    \bottomrule
  \end{tabular}
\end{table}

Without separators, DA improves PER (0.63\,$\to$\,0.58\,\%,
$-8$\,\% relative).
With separators, DA does not improve over CE at 9.7M
(0.52\,$\to$\,0.53\,\%).
Separators alone provide the largest single-factor PER improvement
($-17$\,\%).
The best result (0.48\,\%) requires \emph{all three factors}:
17.2M capacity + DA + separators + distance override.

\textbf{Capacity interaction.}
At 17.2M without separators, DA~Loss does not improve PER relative
to CE (0.61\,\% vs.\ 0.60\,\%).
However, at 17.2M \emph{with} separators and corrected distances,
DA achieves the best PER observed (0.48\,\%).
This refutes the simple hypothesis that ``large models memorize and
DA interferes''---the relevant variable is the interaction between
capacity, structural representation, and distance correction, not
capacity alone.
Dedicated ablation is required for definitive conclusions about
memorization.

\subsection{Comparison with LatPhon}

LatPhon~\cite{chary2025latphon} reports PER~0.86\,\%
(Wilson CI [0.56\,\%, 1.16\,\%], $n\!\approx\!500$).
Our 17.2M DA+dist model achieves PER~0.48\,\%
(CI [0.46\,\%, 0.51\,\%], $n\!=\!28{,}782$).
Both systems report Wilson CIs over the same lexical resource
(ipa-dict); the intervals do not overlap---a statistically
significant difference at 95\,\% confidence.
Methodological differences preclude direct architectural comparison:
LatPhon is a multilingual Transformer; our system is a monolingual
BiLSTM with DA~Loss.
The result is consistent with the hypothesis that methodological
design (loss + evaluation protocol) can compensate for architectural
differences.

% =============================================================
\section{Analysis}
\label{sec:analysis}
% =============================================================

\subsection{Dominant Error Patterns}

Table~\ref{tab:confusions} shows the top confusions.
Over 60\,\% of errors are vowel neutralizations---phonological
ambiguity inherent to PT-BR.
The global /e/:/\textipa{E}/ corpus ratio is 7.1:1, but by position
the pre-tonic ratio is 24.9:1 (/e/~dominant) while the tonic ratio
is 0.33:1 (/\textipa{E}/~dominant).
The model overgeneralizes the pre-tonic /e/ bias to tonic
syllables---a corpus distribution artifact requiring explicit
prosodic features to resolve.

\begin{table}[t]
  \caption{Top phoneme confusions, 17.2M DA+dist}
  \label{tab:confusions}
  \centering\small
  \begin{tabular}{lrl}
    \toprule
    Substitution & Count & Mechanism \\
    \midrule
    \textipa{E}$\to$e & 255 & Unstressed neutralization \\
    e$\to$\textipa{E} & 197 & Same (reverse) \\
    \textipa{O}$\to$o & 131 & Mid-back neutralization \\
    i$\to$e           & 121 & Vowel reduction \\
    o$\to$\textipa{O} & 95  & Mid-back neutralization \\
    \bottomrule
  \end{tabular}
\end{table}

\subsection{Error Quality: DA Loss Redistribution}

Table~\ref{tab:classdist} quantifies the redistribution.
DA~Loss reduces Class~D errors \textbf{19\,\% relative}
(4.3M CE $\to$ 9.7M DA), independently of PER improvement.
In downstream TTS, Class~B errors (e.g.,
/e/$\leftrightarrow$/\textipa{E}/) are perceptually less salient
than Class~D (vowel$\leftrightarrow$stop); formal MOS/ABX validation
remains future work.

\begin{table}[t]
  \caption{Error class distribution. $\dagger$~Class~D inflated by
           structural token confusions.}
  \label{tab:classdist}
  \centering\small
  \begin{tabular}{lcccc}
    \toprule
    System & PER & Cls~B & Cls~D & D/errors \\
    \midrule
    CE base.\ (4.3M)            & 0.66\,\% & 0.39\,\% & 0.54\,\% & 50.9\,\% \\
    DA $\lam\!=\!0.2$ (9.7M)    & 0.58\,\% & 0.36\,\% & 0.44\,\% & 48.4\,\% \\
    DA+dist (17.2M+sep.)        & 0.48\,\% & 0.29\,\% & 0.53\,$\%^\dagger$ & 47.4\,\% \\
    \bottomrule
  \end{tabular}
\end{table}

\subsection{OOV Generalization}
\label{sec:oov}

The 17.2M DA+dist model was evaluated on 31~curated OOV words across
6~categories (Table~\ref{tab:oov}).
The 5/5 result on genuine novel PT-BR words---verified absent from
training---includes correct: palatalization (d+i$\to$d\textipa{Z}),
coda reduction (l$\to$w), rhotacism (rr$\to$x), and nasal vowel
mapping (om$\to$\~{o}).

The $N\!=\!31$ OOV bank is a \textbf{diagnostic probe}, not a
replacement for the main 28{,}782-word test set.
Its purpose is qualitative: demonstrating that specific phonological
rules transfer to unseen words.
A separate set of 35~neologisms (5~categories: anglicisms,
borrowed verbs, native neologisms, proper names, invented words)
provides additional evidence along the same axis.

\textbf{Defined failure modes}: (1)~geminate consonants from
Italian/English loans, absent from training;
(2)~English-phonology anglicisms, genuinely OOV phonologically;
(3)~characters k, w, y---hard character-level OOV.

\begin{table}[t]
  \caption{OOV generalization, 17.2M DA+dist}
  \label{tab:oov}
  \centering\small
  \begin{tabular}{lcc}
    \toprule
    Category & Correct & Phon.\ Score \\
    \midrule
    PT-BR neologisms        & 6/9\,(67\,\%)   & 97\,\% \\
    Geminate consonants     & 1/5\,(20\,\%)   & 81\,\% \\
    Anglicisms (in-vocab)   & 1/5\,(20\,\%)   & 71\,\% \\
    Char-OOV (k/w/y)        & 0/3\,(0\,\%)    & 68\,\% \\
    \textbf{Real PT-BR OOV} & \textbf{5/5\,(100\,\%)} & \textbf{100\,\%} \\
    Controls (in training)  & 4/4\,(100\,\%)  & 100\,\% \\
    \midrule
    \textbf{Total}          & \textbf{17/31\,(55\,\%)} & --- \\
    \bottomrule
  \end{tabular}
\end{table}

% =============================================================
\section{Discussion}
\label{sec:discussion}
% =============================================================

\subsection{Memorization vs.\ Rule Learning}

The 60/10/30 split is deliberate: we prioritize a large test set
over maximizing training data.
With 17.2M parameters and 57{,}562 training words, the model
\emph{could} memorize---a hash table would be more compact.
Yet the 100\,\% accuracy on genuine OOV PT-BR words (correct
palatalization d+i$\to$d\textipa{Z}, coda reduction l$\to$w,
nasal mapping om$\to$\~{o}) is consistent with phonological rule
acquisition.
Furthermore, reducing training data from 60\,\% to 50\,\% (Exp105)
degrades PER only 10\,\% relative (0.49\,$\to$\,0.54\,\%),
suggesting the corpus is well above the saturation threshold for
this architecture---additional training data yields diminishing
returns.

\subsection{PER/WER Trade-off with Separators}

The systematic trade-off between PER improvement ($-17$\,\% to
$-20$\,\%) and WER degradation ($+6$\,\% to $+8$\,\%) with
syllable boundary tokens is structural, not tunable.
Each misplaced separator triggers a full word error under
exact-match WER---the model has more opportunities for error per
output sequence.
This trade-off is independent of architecture, capacity, and loss
function, persisting across all configurations tested.

\textbf{Practical recommendation}: use the 9.7M DA model (no
separators, WER~4.96\,\%) for word-level lookup and NLP tasks;
use the 17.2M DA+dist model (with separators, PER~0.48\,\%) for
TTS and forced alignment where sub-word accuracy matters.

\subsection{Limits of DA Loss}

Two boundaries are identified.
\textbf{Structural tokens}: PanPhon assigns zero vectors to
non-phonemic symbols (``.'', ``\textipa{"}''). The distance
override corrects this partially, but positional confusions
persist---the model distinguishes the tokens but still misplaces
them in the sequence.
\textbf{Signal scale}: DA is bounded by
$\lam\!\times\!1.0\!\times\!1.0 = 0.20$, while early-training
CE reaches $\sim$16.
DA contributes $<$5\,\% of the gradient signal when the model is
far from convergence, becoming effective primarily in the
transition zone (CE~0.3--1.5, epochs 30--80).

\subsection{Future Work}

Three directions are promising:
(1)~\textbf{phonotactic constraints}: PT-BR restricts codas to
4~consonants and $\sim$78\,\% of syllables are open (CV)---an FSA
encoding syllable structure could reduce separator placement errors;
(2)~\textbf{continuous articulatory space}: replacing PanPhon's
24-D binary features with a 7-D continuous vocal tract space would
assign distinct coordinates to structural tokens, eliminating the
need for distance overrides;
(3)~\textbf{morphosyntactic context}: homographic heterophones
(e.g., \textit{jogo} as noun /\textipa{O}/ vs.\ verb /o/) require
sentential context beyond word-level G2P.

% =============================================================
\section{Conclusion}
% =============================================================

We presented Distance-Aware (DA) Loss, a phonologically-graded
training objective for G2P that penalizes substitution errors
proportionally to PanPhon articulatory distance, weighted by
prediction confidence.
Applied to Brazilian Portuguese, DA~Loss reduces catastrophic
(Class~D) substitutions by 19\,\% relative while achieving
PER~0.48\,\% (Wilson CI $\pm$0.03\,pp) and WER~5.33\,\% in the
reference configuration, with a complementary no-separator
configuration reaching WER~4.96\,\%, on the largest PT-BR G2P
evaluation reported---28{,}782 stratified words.
A factorial analysis reveals that the best performance requires
the interaction of three factors---capacity, syllable separators,
and DA~Loss with corrected structural distances---rather than any
single factor.
Genuine OOV generalization (5/5 novel PT-BR words) supports
phonological rule acquisition beyond corpus memorization.

DA~Loss is applicable to any language with PanPhon coverage;
effectiveness beyond PT-BR remains to be validated experimentally.
Limitations: geminate reduction for absent loan words,
English-phonology anglicisms as genuine phonological OOV, and
homograph disambiguation requiring morphosyntactic context beyond
word-level G2P.

% =============================================================
% ACKNOWLEDGMENTS — commented out for double-blind review.
% Restore for camera-ready version.
% =============================================================
% \section*{Acknowledgments}
% [Restored in camera-ready]

% =============================================================
\section*{Use of Generative AI Tools}
% =============================================================

Generative AI tools (Claude, Anthropic) were used during manuscript
preparation exclusively as editorial assistance: grammar and spell
checking, bibliographic reference search, citation management,
terminology consistency verification across sections, and \LaTeX{}
formatting.
The tool did not generate new text, hypotheses, figures, code, or
data interpretations.
All scientific content---experimental design, source code
implementation, experiment execution, data collection and analysis,
interpretation, and conclusions---is the sole work of the authors.

% =============================================================
\bibliographystyle{IEEEtran}
\bibliography{slt_refs}
% =============================================================

\end{document}
